%\begin{quote}
%"Les archives audiovisuelles doivent satisfaire à des normes sévères en ce qui concerne la documentation de leurs acquisitions, des communications et d’autres transactions afin de démontrer le sérieux de leur gestion. Du fait de la complexité de leurs fonds, la précision de l’information est ici essentielle." \footcite{zotero-item-772} 
%\end{quote}
%L'indexation de fonds d'archives doit être "globale pour rendre compte de la finalité du fonds et de sa structure" et "locale pour permettre l’accès et la consultation des unités documentaires composant le fonds"
%"le cinéma a la capacité de traiter n'importe quel fait saisi dans le temps, de prendre tout ce qu'il veut de la vie"

%ou alors : accroche sur le fait que l'indexation manuelle ne suffit plus
%Or, l'indexation = une représentation de la nature des contenus et de leur localisation dans le fonds pour le rendre accessible et exploitable." \footcite{bachimont}
%"Il est nécessaire d’avoir une représentation de la nature des contenus et de leur localisation dans le fonds pour le rendre accessible et exploitable." \footcite{bachimont2007} 
%indexation : permise par un système de gestion des collections 


\begin{quote}
	La nature audiovisuelle des documents [...] leur confère le statut d’objets temporels. Il se pose alors le problème spécifique à l’audiovisuel de parcours et de navigation dans un objet temporel [...] pour le rendre manipulable dans le cadre d’une exploitation donnée \footcite{bachimont}. 
\end{quote}

Dès 1998, le chercheur Bruno Bachimont exprime la manière dont la manipulation des contenus audiovisuels, au sein des pratiques documentaires, est contrainte par leur nature temporelle, fixée sur un support. En effet, les documents audiovisuels comprennent "des images et / ou des sons reproductibles réunis sur un support matériel dont : l’enregistrement, la transmission, la perception et la compréhension exigent le recours à un dispositif technique" \footcite{zotero-item-772}. Le suport, numérique ou analogique, contraint la consultation du contenu, poussant les institutions patrimoniales à privilégier une "consultation indirecte" via sa description \footcite{bachimont}. \\
De fait, l'appui sur un logiciel de gestion des collections joue un rôle déterminant dans la mise en valeur et l'organisation de la description des fonds audiovisuels. Dispositif avant tout technique, il permet de gérer les "informations structurées" produites par les tâches d'indexation, conformément à des règles métiers, en fournissant notamment un modèle de données, afin d'organiser la description en couches intelligibles et interreliées \footcite{zotero-item-766}. \\
Or, la structuration de la description  au sein d'un tel outil (indexation) présuppose justement l'analyse approfondie du contenu audiovisuel comme "objet temporel", via l'étude de sa syntaxe en différents plans et séquences, permettant de concrétiser une "navigation" précise \footcite{bachimont}. Cette analyse précise dont un tel type de document devrait être l'objet, est toutefois contrainte par la densification des collections patrimoniales \footnote{Cette densification est produite par la nature même du document audiovisuel, dont la reproductibilité technique permet la copie multiple à l'identique des exemplaires.}. Il devient ainsi difficile de traiter l'ensemble des documents audiovisuels, via leur lecture sur un support, leur analyse fine et la structuration de leur description. 
Si le logiciel permet de structurer la description produite, il ne permet toutefois pas encore de palier à ces contraintes de temps et de moyens. \\
Or justement, le développement croissant de l'intelligence artificielle offre des fonctionnalités innovantes afin d'automatiser certaines pratiques documentaires, dont l'indexation. L'automatisation de l'indexation appliquée aux collections audiovisuelles permettrait l'analyse et la description du flux audiovisuels de manière efficace, afin d'y donner largement accès. L'éditeur de gestion des collections, dont la logique commerciale encourage l'innovation, possèderait alors rôle pivot dans la proposition de fonctionnalités automatiques aux intitutions, permettant ainsi d'accélerer et d'améliorer la découvrabilité de documents complexes. \\
Néanmoins, si certaines tâches classiques d'analyse du contenu de l'image sont démocratisées, y compris en milieu patrimonial, la compréhension fine de la syntaxe audiovisuelle reste encore sous développée, puisqu'elle nécessite a priori des algorithmes lourds et couteux. C'est notamment le cas de la segmentation du flux audiovisuel en les différentes séquences qui le composent, qui permettrait pourtant une description fine du contenu afin d'établir une "navigation" au sein du document \footcite{bachimont}. \\
Une double complexité se pose ainsi. D'une part, la nécesité de devélopper une automatisation de l'indexation des documents audiovisuels est contrainte par l'inexploitation d'une compréhension fine de leur syntaxe en différentes séquences, par les algorithmes d'intelligence artificielle existants. D'autre part, l'implémentation de processus automatisés au sein d'un outil de gestion densifie la production de l'information, complexifiant de fait la structuration du modèle de données exitante, propre à des normes et règles métier contextuelles. 
%pour les documents textuels ou même iconographiques, l'émergence des modèles fondés sur l'intelligence artificielle, offre des possibilités de description automatique efficaces, certaines tâches appliquées à l'audiovisuel restent encore sous développées. C'est notamment le cas de la segmentation d'un document audiovisuel en les différentes séquences qui le composent, qui permettrait une description fine du contenu et d'établir une "navigation" \footcite{bachimont}. Cette tâche implique une compréhension fine du langage temporel propre à l'audiovisuel et nécessiterait a priori des algorithmes lourds et couteux. 



%Pourtant, une réelle nécessité subsiste concernant ce type d'archives puisque, nécessitant du matériel de lecture, on privilégie la "consultation indirecte", via l'indexation de description, plutot qu'une "consultation directe", nécessitant du matériel spécifique \footcite{bachimont}. 
%Ainsi, contrairement aux documents imprimés et textuels dont la forme rend accessible le contenu, les supports audiovisuels sont opaques et rendent la description documentaire fastidieuse, et d'autant plus leur indexation, qui rend la description accessible et exploitable par des utilisateurs \footcite{bachimont}. 
%Ainsi, ces "objets temporels" nécessitent un "dispositif technique", qu'il soit numérique ou analogique, capable d'en réaliser la lecture et d'en permettre la manipulation \footcites{bachimont}{zotero-item-772}. 
%(face à la masse des documents : décrire de manière précise des documents sélectionnés : séquences, plutôt que de tout décrire grossièrement) Or, pour éviter une consultation directe, la description doit détailler le caractère temporel du contenu et sa syntaxe, notamment constitué d'unités telles que les plans ou les séquences,  L'indexation manuelle ne peut pas atteindre ce niveau de précision, qui supposerait un temps dédié considérable. 

C'est précisément cette double difficulté qui anime l'étude de ce mémoire de recherche, réalisé dans le cadre d'un stage de six mois au sein de l'entreprise Skinsoft, éditrice de solutions logicielles de gestion des collections. Effectué au sein de l'équipe métier, initiatrice de fonctionnalités innovantes adaptées aux besoins du patrimoine, ce stage se porte spécifiquement sur l'amélioration générale d'un espace de travail dédié à la gestion des archives audiovisuelles (S-movies), fondé sur le FRBR. L'amélioration de cette interface est motivée par le projet d'automatiser certaines tâches d'indexation à l'aide de l'intelligence artificielle. Ces tâches, issues d'une demande client spécifique, comprennent la segmentation automatique de contenus audiovisuels en séquences, puis leur indexation. Les réflexions et observations, à la fois techniques et métier, qui ont traversé notre stage, constituent le point de départ de ce mémoire, qui examine les conditions de faisabilité et les modalités d'intégration d'une indexation automatique appliquée à l'audiovisuel au sein d'une modélisation adaptative du FRBR. Si le projet n'a pas été développé en cours de stage, notre étude examine toutefois les implications et les enjeux soulevés par la mise en place d'une automatisation de l'indexation des collections audiovisuelles, en accord avec les besoins métier existants.

Dès lors, dans quelle mesure la variabilité du modèle de données FRBR adapté aux archives audiovisuelles conditionne-t-elle la mise en place d'une segmentation, et d'une indexation automatiques pour ce type d'archives ? Quels types d'algorithmes, fondés sur l'intelligence artificielle, peuvent être proposés au sein d'un outil de gestion des collections, à destination d'institutions patrimoniales ?

Dans une première partie, nous analyserons les fondements du modèle de données FRBR tel qu'il a été pensé pour le milieu des bibliothèques, avant de nous pencher sur les modalités de son adaptation aux archives et documents audiovisuels. Nous prendrons pour exemple l'espace de travail S-Movies, développé par Skinsoft, à travers l'étude d'un corpus audiovisuel spécifique. La présentation contextuelle du modèle de données nous amènera à examiner les complexités documentaires posées par le projet de la mise en place d'une segmentation et d'une indexation automatiques des documents audiovisuels. \\
Dans une seconde partie nous nous focaliserons d'une part sur les technologies existantes pour l'automatisation du traitement audiovisuel, avant de nous pencher sur les outils disponibles pour la segmentation automatique du flux audiovisuel en séquences. Nous proposerons une mise en application de la segmentation à travers l'outil pyscenedetect, en mettant à l'épreuve le corpus audiovisuel présenté en première partie. D'autre part, à partir du même corpus, nous proposerons la description des segments produits par mots-clés, à l'aide du modèle YOLO, dont nous examinerons les modalités d'entraînement et de validation. Nous nous pencherons finalement sur l'indexation des données produites automatiquement au sein du modèle de données FRBR, à travers une validation utilisateur, afin d'améliorer la découvrabilité des documents.  \\
Dans une troisième partie, nous examinerons la manière dont le contrôle de la qualité des données produites pourra s'effectuer au sein d'un logiciel de gestion des collections, afin d'intégrer l'utilisateur à la boucle automatique. Nous proposerons par ailleurs des perspectives d'améliorations d'interface afin d'assurer la qualité des données et d'intégrer l'automatisation. Enfin, nous évaluerons la mesure dans laquelle un tel processus pourra être généralisé à d'autres institutions, au regard des besoins existants concernant l'indexation et la segmentation.  

%C'est pourquoi face à l'augmentation croissante des productions de documents et des acquisitions patrimoniales, la mise en place de solutions automatiques au sein d'outils de gestion permettrait d'accélérer la structuration de l'information afin de pouvoir y donner largement accès. \footcite{barreaux2017}
%En effet, la temporalité des documents audiovisuels limite la description "locale" et précise de leur contenu, qu'il faudrait segmenter en unités temporelles réduites afin d'en préparer l'indexation \footcite{bachimont}.
%Mise en place d'une segmentation automatique : 1er pb : comment faire comprendre à la machine la temporalité audiovisuelle ?
%
%
%7. indexation : pratique normée : modèle de données 
%deuxième pb : comment intégrer ces fonctionnalités automatiques dans un logiciel de gestion des collections dont la structure est conditionnée par des normes et pratiques documentaires normées. 
%
%
%
%Intégrer un état de l'art : sur l'automatisation de l'indexation des archives audiovisuelles : combinaison de plusieurs algorithmes, pas d'outils concret qui feraient l'ensemble du processus d'indexation. Ce mémoire tente d'articuler un pipeline d'indexation automatique avec un modèle FRBR adaptatif  
%
%
%
%
%
%
%OLD 
%
%Dans la Philosophie de l'archivistique audiovisuelle publiée par l'UNESCO, Edmonson Ray pose les fondements de la complexité de la gestion et de l'indexation des documents audiovisuels en raison de la nature de leur support, une indexation qui doit alors être soumise à une structure définie et claire pour en instaurer la compréhension. L'auteur définit les documents audiovisuels comme suit : 
%
%\begin{quote}
%"Constituent des documents audiovisuels les œuvres comprenant des images et / ou des sons reproductibles réunis sur un support matériel dont : l’enregistrement, la transmission, la perception et la compréhension exigent le recours à un dispositif technique;le contenu visuel présente une durée linéaire"
%\footcite{zotero-item-772}
%\end{quote}
%
%
%% : préparation et gestion de l’entrée dans les collections, prise d’inventaire, documentation, gestion des informations liées à la conservation, récolement, post-récolement, gestion des multimédias, des droits d’auteur, des mouvements, mais aussi diffusion au public (sous forme d'extraction de données ou de catalogue numérique), outil de médiation, etc 
%
%%La collection incarne, selon la même source, un "ensemble non fini d'objets matériels ou immatériels, réunis et classés par un individu ou une institution, en raison de leur valeur scientifique, artistique, esthétique, documentaire, affective, pour leur prix, leur rareté, etc. \footcite{zotero-item-766} 
%
%Ainsi défini, un outil de gestion des collections permet la réalisation des diverses actions liées au cycle de vie de l'ensemble des objets de collection préservés au sein d'une instutition.
%
%La réalisation de ces actions repose nécessairement sur la production "d'informations structurées" exploitables au sein du système. Cette structuration est permise par l'indexation, « l’opération qui consiste à décrire et à caractériser un document à l’aide de représentations des concepts évoqués dans ce document ». \footcite{zotero-item-767} 
%
%L'indexation permet précisément d'intelligibiliser un objet au sein d'un système de gestion des collections et de le rendre repérable et interrogeable auprès de l'utilisateur. Dans le cas des collections audiovisuelles, la description d'un objet existe à plusieurs niveaux : contenu intellectuel de l'oeuvre, versions de l'oeuvre, publications de l'oeuvre, localisation et état matériel de ses exemplaires... Une indexation simple ne permet pas la hiérarchisation de ces différents niveaux de description, or, ils facilitent la gestion de l'objet. La mise en place d'une modélisation des données d'indexation permet notamment de les structurer et de les hiérarchiser, afin de mieux appréhender la complexité de la description de tels contenus. 
%
%Un outil de gestion des collections peut ainsi s'appuyer sur un modèle de données, tel que le FRBR (\textit{Functional Requirements for Bibliographic Records}), modèle que nous analyserons au cours de cette étude. Issu du milieu des bibliothèques, il est également largement utilisé pour la gestion de collections audiovisuelles dont la nature se prête particulièrement à une indexation à différents niveaux. En effet, un document audiovisuel, reproductible par nature, se caractérise par son contenu intelectuel et ses multiples versions et publications, ainsi que ses exemplaires physiques. Ce modèle, implanté au sein d'un outil de gestion des collections, permet alors de dissocier le contenu intellectuel d'un objet de son support matériel. Le FRBR est toutefois variable selon la collection audiovisuelle, dont la nature conditionne les niveaux d'indexation nécessaires. Une collection expérimentale non diffusée ne nécessite pas les mêmes niveaux d'indexation qu'une collection destinée au grand public. Ainsi, la structure de l'indexation des collections au sein d'un outil de gestion peut être adapté en fonction de la nature de la collection gérée au sein de chaque institution.
%
%Par ailleurs, si la modélisation des données d'indexation permet d'adapter l'outil de gestion des collections à la complexité descriptive des collections audiovisuelles, le volume de ces collections ne permet pas à une institution de rendre compte de la globalité de ses objets au sein d'un tel outil. C'est pourquoi face à l'augmentation croissante des productions de documents et des acquisitions patrimoniales, la mise en place de solutions automatiques au sein d'outils de gestion permettrait d'accélérer la structuration de l'information afin de pouvoir y donner largement accès. \footcite{barreaux2017}
%
%
%Les tâches d'indexation permises par un outil de gestion des collections restent en effet répétitives et chronophage, et font désormais l'objet de tentatives  d'automatisation à l'aide de l'intelligence artificielle. Dans le cas de documents textuels ou iconographiques, l'émergence des modèles fondés sur l'intelligence artificielle, offre des possibilités d'extraction automatique de métadonnées efficaces sur des documents numérisés. Toutefois, les collections audiovisuelles, résistent à la mise en place d'une extraction automatique de leur contenu, en raison de leur format "boite noire" \footcite{zotero-item-582}.
%
%Ainsi, d'une part, l'indexation automatique de ces fonds supposerait que les modèles d'intelligence artificielle soient capables de lire et de comprendre le langage cinématographique, caractérisé par la mise en mouvement d'images fixes, le son, et le montage. Or, la compréhension automatique de ce langage reste inégale, en particulier lorsqu'elle est confrontée à des fonds d'archives audiovisuels patrimoniaux, anciens ou dégradés, dont la qualité sonore et visuelle est compromise. A priori, l'étude d'un tel type de documents nécessiterait des algorithmes d'intelligence artificielle lourds et couteux, incompatibles avec les logiques institutionnelles patrimoniales. Le rôle de recherche et de proposition d'une indexation automatique pour les collections audiovisuelles reviendrait à l'éditeur de l'outil de gestion des collections, dont la logique produit encourage l'innovation constante. 
%Toutefois cette innovation se heurte à une double difficulté que nous avons posée : la complexité pour les modèles automatiques de lire et de comprendre un contenu visuel, sonore et temporel, ainsi que l'inscription de fonctionnalités automatiques au sein d'un modèle de données (FRBR) variable, hautement dépendant du contexte institutionnel. 
%C'est précisément cette double difficulté qui anime l'étude de ce mémoire de recherche, réalisé dans le cadre d'un stage de six mois au sein de l'entreprise Skinsoft, éditrice de solutions logicielles de gestion des collections. Effectué au sein de l'équipe métier, forte de propositions de  fonctionnalités innovantes adaptées aux besoins des clients institutionnels, ce stage se porte spécifiquement sur l'amélioration de l'espace de travail dédié aux archive audiovisuelles (S-movies), fondé sur le FRBR, à travers le projet d'automatiser certains pans de l'indexation à l'aide de l'intelligence artificielle. Les réflexions et observations d'un point de vue métier qui ont traversé ce stage constituent le point de départ de ce mémoire, dont l'étude examine les conditions de faisabilité et les modalités d'intégration d'une indexation automatique appliquée à l'audiovisuel au sein d'un système de gestion de collections fondé sur une modélisation adaptative du FRBR.
%
%Intégrer un état de l'art : sur l'automatisation de l'indexation des archives audiovisuelles : combinaison de plusieurs algorithmes, pas d'outils concret qui feraient l'ensemble du processus d'indexation. Ce mémoire tente d'articuler un pipeline d'indexation automatique avec un modèle FRBR adaptatif  
%
%Dans quelle mesure la variabilité du modèle de données FRBR adapté aux archives audiovisuelles conditionne-t-elle la mise en place d'une indexation automatique pour ce type d'archives ? Quels types d'algorithmes fondés sur l'intelligence artificielle peuvent être proposés au sein d'un outil de gestion des collections adaptatif à la nature des collections audiovisuelles ?
%Dans une première partie, nous analyserons les fondements du modèle FRBR et les modalités de son adaptation aux archives audiovisuelles, afin d'examiner les questions documentaires posées par le projet de la mise en place d'une indexation automatique de tels documents. 
%Dans une seconde partie nous étudierons la possibilité d'une implémentation concrète d'outils d'indexation automatique au sein d'une modélisation FRBR d'un outil de gestion des collections. Nous y analyserons la faisabilité technique d'une mise en place de l'indexation à partir d'un séquençage automatique de fichiers numériques audiovisuels et d'une extraction automatique de mots-clés, pour en analyser le reversement au sein des différents niveaux de description modélisés par le FRBR.
%Enfin, nous évaluerons dans une troisième partie 
%les limites, biais et enjeux éthiques liés à l'inadaptation des données d'entraînement d'IA aux fonds d'archives. Nous montrerons alors que les limites du « tout-automatique » invitent à réorienter la conception des logiciels métier vers un paradigme d'**indexation augmentée et d'intelligence hybride**, plaçant l'interaction homme-machine au cœur de la gouvernance de la qualité documentaire.
%
%
%%Le FRBR, issu du milieu des bibliothèques et élaboré par l'IFLA en 1998, puis révisé sous le modèle LRM, est fondé sur une approche entité-relation au sein de laquelle la description de l'objet documentaire en répartie en quatre niveaux : l' \textit{Œuvre} (la création intellectuelle), l' \textit{Expression} (sa réalisation linguistique ou artistique), la \textit{Manifestation} (sa matérialisation éditoriale) et l' \textit{Item} (l'exemplaire physique ou numérique conservé). Ce modèle permet de dissocier le contenu intellectuel d'un objet de son support matériel, dissociation que l'outil de gestion des collections doit concrétiser. Il est adapté, au sein d'un outil de gestion des collections, en fonction de la nature des collections au sein de chaque institution. En effet, si la modélisation FRBR est pensée en premier lieu pour la description d'ouvrage bibliographiques, le modèle est également largement utilisé pour la description des archives et en particulier celle des archives audiovisuelles. Le modèle est alors réadapté en fonction de la nature des objets décrits afin de proposer un mode de description et une structure de données pertinents. Un outil de gestion des collections peut alors proposer une modélisation adaptative à raison des collections de nature multiples préservées par les institutions patrimoniales. 

%D'autre part, la modélisation des données audiovisuelles pose un défi théorique tout aussi complexe. Pour s'adapter aux spécificités du médium cinématographique, la Fédération Internationale des Archives du Film (FIAF) a révisé le modèle FRBR original en fusionnant les niveaux \textit{Œuvre} et \textit{Expression} et en introduisant la notion de \textit{Variante} (passant d'une structure WEMI à une structure WVMI). En effet, une œuvre audiovisuelle est indissociable de son expression, mais elle existe à travers une multitude de variantes (montages alternatifs, censures, doublages, restaurations). Si la FIAF propose un cadre théorique, le FRBR reste un modèle extrêmement flexible, réinterprété par chaque institution en fonction de la nature de ses corpus et de ses règles métiers propres.




 

%\chapter{Partie 1 - Indexer les archives audiovisuelles}  
%
%\begin{itemize}
%
%\section{Le FRBR, modèle d’indexation conçu pour les bibliothèques}
%
%
%
%\item {Un modèle entité-relation qui permet la description de concepts reliés entre eux. Avant ce modèle, une notice bibliographique correspondait à un objet physique. Il permet par ailleurs de mieux décrire les objets complexes comme les archives audiovisuelles.}  
%
%\item{La structure WEMI, qui permet de décrire chaque niveau : Oeuvre, expression, manifestation, item. C’est une structure pensée pour des oeuvres textuelles, avec une verticalité limitée }
%
%\item{Les objectifs du FRBR sont les besoins utilisateurs. Il doit pouvoir répondre à quatre tâches utilisateurs : trouver, identifier, sélectionner, obtenir. Afin de remplir ces objectifs, une interface intuitive doit pouvoir accueillir le modèle FRBR.  
%}
%
%\end{itemize}
%
%
%\section{Adapter le FRBR aux archives audiovisuelles}
%
%\begin{itemize}
%
%
%\item{Remodéliser le FRBR sur les caractéristiques complexes du médium cinématographique. L’exemple de S-Movies : un espace de travail pour la gestion des archives audiovisuelles, fondé sur une remodélisation du FRBR, adaptable selon les corpus de données} 
%
%\item{L'inscription de différents types de corpus audiovisuels dans ce modèle repensé du FRBR : la cause d’une expérience utilisateur complexe et non-intuitive }
%
%\item{Le projet d’ajouter sur cette remodélisation du FRBR, une indexation augmentée et automatisée des collections audiovisuelles}
%	
%\end{itemize}
%
%\section{Automatiser l’indexation des collections audiovisuelles basée sur le FRBR} 
%
%\begin{itemize}
%
%
%\item{L’automatisation de l’indexation des archives audiovisuelles, motivée par les caractéristiques physiques de ce type de collection.  En effet, les collections audiovisuelles sont indexées par leur contenu puisque leur support physique est qualifié de “boite noire”. Leur indexation est donc un travail long et fastidieux. } 
%
%\item{Les possibilités d’automatisation de l’indexation sur les collections audiovisuelles : 
%
%Analyse d’objets, (segmentation sémantique) 
%
%Analyse et transcription de texte  
%
%Analyse et tanscription audio  
%
%Reconnaissance vocale et reconnaissance des locuteurs 
%
%Extraction d’entités nommées (lier ces analyses automatiques à des référentiels) }
%
%\item{Séquençage (segmentation temporelle)  : une fonctionnalité spécifique aux archives audiovisuelles puisqu’il s’agit de comprendre les relations entre les plans pour pouvoir adéquatement les séquencer.  
%
%Faire de l’adaptation de cette automatisation au modèle FRBR dédiée aux archives audiovisuelles, un objectif } 
%
%
%\end{itemize}
%
%\chapter{Partie 2 -  Implémenter l’indexation automatique des archives audiovisuelles au sein d’une remodélisation du FRBR} 
%
%
%
%\section{Le projet d’implémentation de l’IA chez Skinsoft (projet Skai), dont l'analyse vidéo est une composante majeure} 
%
%\begin{itemize}
%
%\item{L’objectif du projet est de proposer différents moteurs IA spécialisés dans différents tâches automatisées, entraînés sur des données ouvertes patrimoniales (POP, Joconde, Europeana). Cela permettrait aux modèles d'être adaptés au contexte patrimonial et de produire des métadonnées conformes aux normes bibliographiques. En revanche, ces bases de données contiennent majoritairement des archives textuelles et iconographiques. }
%
%\item{L’orientation donnée est de proposer des modèles force de proposition, et facilitateurs pour l’utilisateur : l’utilisateur doit pouvoir garder le contrôle sur ses données }
%
%\item{Parmi les moteurs IA en cours de développement : un moteur spécialisé pour l’analyse vidéo multimodale, adapté aux collections audiovisuelles. Une des composantes principales de ce moteur, issue d’une demande client, est la segmentation temporelle, c’est-à-dire le séquençage vidéo. L’objectif est, à partir des séquences produites, d’apposer des mots clés sur le contenu de l’image.  }
%
%
%\end{itemize}
%
%
%
%\section{Séquençage et analyse d’un corpus audiovisuel (exemple du corpus de musée Albert Kahn)}
%
%\begin{itemize}
%
%\item{Séquençage des archives audiovisuelles en unités analysables, la première étape complexe d’un pipeline multimodal, dépendante du type de corpus analysé : la détection des coupures permet une classification des segments uis une vérification humaine.}  
%
%\item{Combinaison des approches syntaxique et sémantique pour la segmentation de scènes : réseaux de neurones récurrents (RNN), réseaux de neurones convolutifs (CNN) ou Transformers ? }
%
%\item{L’adaptation de ces outils au corpus Albert Kahn, archives audiovisuelles des années 1920, documentaires et muettes : un séquençage facilement identifiable mais difficilement classifiable d’un point de vue sémantique.} 
%
%\end{itemize}
%
%
%\section{Reversement des métadonnées dans le modèle FRBR : dans quelle mesure les résultats produits peuvent-ils s’inscrire au sein de ce modèle d’indexation ? } 
%
%\begin{itemize}
%
%
%\item{Cartographie du reversement des métadonnées : quels résultats pour quels niveaux du FRBR ? Cette cartographie dépend de l’utilisation du FRBR par l’institution concernée. }
%
%\item{Le reversement adéquat des métadonnées produites sur différents niveaux du FRBR ne peut se produire sans une interface navigable et intuitive entre ces différents niveaux. Une refonte est ainsi nécessaire pour la mise en place d’un reversement automatique. Actuellement, il existe une confusion des liens entre les niveaux, qui fait l’objet d’un projet de refonte dans l’espace de travail S-Movies. } 
%
%\item{La mise en place d’une interface dans laquelle l’utilisateur pourrait valider les données produites, et décider de leur reversement de tel ou tel niveau du FRBR. Le moteur d’IA aurait alors bien ce rôle de “facilitateur” en proposant des solutions, tout en laissant le choix final à l’utilisateur. Cette solution pourrait rendre la compréhenion du FRBR plus claire. } 
%
%
%\end{itemize}
%
%
%\chapter{Partie 3 - Biais, limites et enjeux }
%
%
%
%\section{Les archives audiovisuelles comme données d’entraînement} 
%
%\begin{itemize}
%
%\item{La sous-représentation des archives audiovisuelles anciennes dans les corpus d’entraînement. Il faudrait entraîner le modèle sur les corpus individuels de chaque institution pour pouvoir produire des résultats pertinents. L’inadaptation des corpus d’entraînement entraîne des biais d’indexation.}  
%
%\item{le manque de temps et d’argent dédié à l’entraînement d’un modèle sur de multiples corpus. Quelles solutions ? Fine tuning, transfert learning, données synthétiques. }
%
%\item{Trouver le juste milieu entre un pipeline complexe multimodal qui demande beaucoup de puissance de calcul, et un pipeline léger mais moins performant. Ce choix dépend effectivement du type d’archives à indexer. }
%
%
%\end{itemize}
%
%\section{La difficulté d’adaptation au modèle de données, même avec un entraînement adéquat  }
%
%\begin{itemize}
%
%
%\item{L’utilisation des niveaux du FRBR varie pour chaque institution, puisqu’elle dépend de la nature du corpus audiovisuel. Le modèle est donc interprété par des règles métiers internes et propres à chaque institution. Il faudrait ainsi entraîner un modèle IA sur ces règles métier liées à l’utilisation du FRBR. Par exemple, la collection du musée Albert Kahn s’adapte mal au FRBR puisqu’il s’agit d’archives documentaires, sans réelle diffusion. Un reversement automatique suppose que le modèle d’IA soit entraîné sur le contexte spécifique de chaque institution. }
%
%\item{S’impose alors la mise en place d'une indexation semi automatisée où l’utilisateur choisit, décide pour respecter la cohérence de l’usage du modèle de données. }
%
%
%\end{itemize}
%
%\section{Vers une indexation augmentée plutôt qu’automatisée }
%
%\begin{itemize}
%
%\item{L’interactivité et le dialogue entre le modèle et l’utilisateur comme réponse aux biais et aux manques des données d’entraînement. Il faudrait rendre le processus transparent comme le préconise Skinsoft dans son projet Skai, en rendant accessible les données sur lequelles le modèle a été entraîné, la réflexion effectuée... }
%
%\item{La transparence comme éthique de la transmission du savoir au public : la diffusion des connaissances est facilitée et augmentée par le modèle d’IA. La mise en place d’une interface de dialogue entre le modèle et l’utilisateur en est la condition essentielle }
%
%\end{itemize}	



